\subsection{LEAST UPPER BOUND OF A SUBSET OF R}

THEOREM 3. Every non-empty subset of the real line which is bounded above (resp. bounded below) has a least upper bound (resp. greatest lower bound).

Let A be a non-empty subset of R, bounded above, and let b be an upper bound of A, so that \(A \subset ] \leftarrow, b]\) . For each \(x \in A\) consider the set \(A_x\) of numbers \(\geqslant x\) which belong to A; the sets \(A_x\) form a filter base B on R, since \(A_y \subset A_x\) if \(y \geqslant x\) . Let a be a point of A. For each \(x \geqslant a\) belonging to A, \(A_x\) is contained in the compact interval {[}a, b{]} and thus the filter base B has a cluster point c.~Since the intervals \([x, \rightarrow[\) are closed, c belongs to their intersection and therefore c is an upper bound of A. But, on the other hand, every upper bound z of A is \(\geqslant c\) , otherwise the neighbourhood \(]z, \rightarrow[\) of c would not contain any point of A. Hence c is the least upper bound of A.

We can argue similarly for a non-empty set B bounded below, or else simply remark that -B is non-empty and bounded above, and that if c is the least upper bound of -B, then -c is the greatest lower bound of B.

The least upper bound \(c\) of \(A\) can be characterized by the following two properties:

\begin{enumerate}
\def\labelenumi{(\roman{enumi})}
\item
  For each \(x \in A, x \leqslant c_*\) .
\item
  For each \(a < c\) , there exists \(x \in \mathbf{A}\) such that \(a < x \leqslant c\) .
\end{enumerate}

The least upper bound of a closed set (non-empty and bounded above) belongs to the set and is its greatest element; and the least upper bound of any non-empty subset A of R which is bounded above may be defined as the largest real number in the closure of A.

